% 第11章 交易环节的横向整合
\section{交易环节的横向整合：从预测到仓位、执行与风控}\label{sec:integration}

前面按方法族竖着切了一遍，这一章横着切：沿一条完整策略的生命周期，看每个环节上 AI 做了什么、环节之间又怎么互相牵制。写这一章的理由是，\textbf{落地失败常常不发生在上游的预测，而发生在预测到仓位、仓位到成交、成交到风控这几段衔接处。}

\subsection{组合构建：把预测翻译成权重}

一个微弱的收益预测，要靠组合优化放大成可观的夏普，就得处理均值--方差框架里那个老问题：$\mu$ 的估计误差被 $\Sigma^{-1}$ 放大（“误差最大化”）。机器学习在这一环的贡献有三条路径。

第一是\textbf{更稳健的协方差与收缩估计}。Ledoit 与 Wolf（2004）\cite{ledoitwolf2004} 的收缩估计把样本协方差向结构化目标收缩，显著降低高维下的估计噪声，这几乎是所有现代组合管线的标配。

第二是\textbf{端到端学习组合权重}。Zhang、Zohren 与 Roberts（2020）\cite{zhang2020dlpo} 直接用深度网络输出权重并优化组合效用，Costa 与 Iyengar（2022）\cite{costa2022} 提出端到端风险预算组合。端到端的吸引力在于它不再“先预测、再优化”两段分家，而把交易成本与风险预算直接写进目标；代价是可解释性下降、对训练数据的依赖加深。

第三是\textbf{把约束写进学习}。现实里权重受再多约束——不卖空、行业上限、单票上限、换手限制——两段式做法无法保证预测最优能转成可行权重。约束感知的学习（在损失里加惩罚、或直接参数化可行域）是这一环的前沿。

这里有一条容易被忽视的量化关系：\textbf{组合的换手率由预测的“信噪比”和优化器的“锐度”共同决定}。预测噪声大时，最优权重对噪声极度敏感，换手率飙升；很多“理论夏普很高”的策略，成本后归零，原因就在这。第 \ref{sec:real-empirical} 节给出了这条关系的一组实测对照：同一批特征里，年化换手 14 倍的候选在成本后仍有夏普 1.01，而 182 倍的候选被压到 $-1.49$——换手相差一个数量级，结局相反。

\subsection{一个成本分解的示意计算}

为了把“摩擦如何吃掉 alpha”说得具体，这里给一个\textbf{示意性}计算——所有假设都明示，它是算术演示，不是实证结果。

设日度截面信号 IC$=0.05$，个股日收益的截面标准差 $\sigma_{cs}=2\%$，组合转移系数 $\tau\approx0.5$，则组合的期望毛收益约为
\begin{equation}
  \mathbb{E}[r_{\text{gross}}/\text{日}] \approx \text{IC}\times\sigma_{cs}\times\tau = 0.05\times0.02\times0.5 = 5\text{bp}.
\end{equation}
成本按固定 $5$bp 加平方根冲击计算，冲击率 $=0.02\sqrt{\text{参与率}}$；取参与率 $1\%$（对应 ADV $10^9$、单票约 $10^7$），得冲击率 $20$bp，于是每次调仓的成本率 $=25\text{bp}\times$ 单边换手率。三种调仓频率下：

留意这里 ADV $=10^9$ 只是\textbf{示意假设}，不是普适值。ADV 该取什么，取决于交易的标的：第 \ref{sec:real-empirical} 节在 A 股样本上取逐个股成交额中位数，其中位数为 $7.46\times10^7$ 元，比这个示意值小一个量级；参与率随之抬高，冲击成本按平方根律放大。用常数 ADV 套所有标的，是最容易把成本算轻的一种做法。

\begin{table}[htbp]
\centering
\caption{同一信号在不同调仓频率下的净收益（示意计算）}
\label{tab:costsplit}
\begin{tabular}{@{}lcccc@{}}
\toprule
调仓频率 & 单边换手率 & 持有期毛收益 & 持有期成本 & 持有期净收益 \\
\midrule
每日 & 40\% & 5.0bp & 10.0bp & $-5.0$bp \\
每周 & 15\% & 25.0bp & 3.8bp & $+21.2$bp \\
每月 & 6\% & 105.0bp & 1.5bp & $+103.5$bp \\
\bottomrule
\end{tabular}
\end{table}

\noindent 需说明：表中假设信息衰减慢，因而毛收益随持有期近似线性增长；若信号半衰期很短（例如两三天），长持有期的毛收益会被显著高估，结论需重算。即便如此，方向性结论是稳的——\textbf{同一个信号，在日频下净收益为负，降频到周频以上才转正。}这说明“信号好不好”与“频率选得对不对”是两回事。AI 最常犯的错，恰恰是把一个只有周频价值的信号，放到日频上去交易，再用样本内的零成本回测证明它“有效”。

\subsection{交易执行：摩擦的管理学}

执行环节的目标只有一个：把想交易的量，以尽可能低的总成本变成成交。总成本是时机风险与冲击成本的折中（Almgren--Chriss\cite{almgrenchriss2000}）。这一环上 AI 的作用是把解析解换成数据驱动的动态策略，落地要点有三：

\begin{itemize}
  \item \textbf{冲击模型必须内生化}。市场冲击大致服从平方根律（Grinold \& Kahn 的工程经验与 Almgren 等 2005\cite{almgren2005} 的直接估计都支持这一形态），策略若假设线性冲击，会系统性高估大单的可执行性。
  \item \textbf{用真实的订单簿数据训练}。逐笔成交与订单簿（第 \ref{sec:altdata} 章）是执行策略唯一可信的训练场。
  \item \textbf{把执行嵌入策略而非事后附加}。最优做法是让信号直接知道“交易它会花多少成本”，否则会生成大量“看起来能赚、扣完成本为负”的信号。
\end{itemize}

\subsection{风险管理：AI 的另一条战线}

风险管理不比信号生成简单，而且它更依赖对“尾部”的刻画，而尾部恰恰是样本最少的地方。

\textbf{波动率与相关性建模}。GARCH 一族（Bollerslev 1986\cite{bollerslev1986}）与现代已实现波动率模型是基础；机器学习改进的是对波动率与相关性的动态预测，以及在高维下的稳健估计。

\textbf{尾部风险}。CVaR（Rockafellar \& Uryasev 2000\cite{rockafellar2000}）与极值理论（McNeil \& Frey 2000\cite{mcneil2000}）给出了尾部风险的度量与估计框架；机器学习（分位数回归、深度分位数网络）让条件尾部的预测能吸收更多特征。这一环的困难在于\textbf{尾部事件样本极少，任何在此处的过拟合都极其危险}。

\textbf{回撤控制与波动率目标化}。Chekhlov、Uryasev 与 Zabarankin（2005）\cite{chekhlov2005} 提出回撤约束的优化；Harvey 等（2018）\cite{harvey2018} 系统论证了波动率目标化对组合的影响。监督式/强化式的回撤控制是近年的活跃方向，但其样本外稳健性同样受制于尾部样本稀缺。

\textbf{制度识别}。Hamilton（1989）\cite{hamilton1989} 的 regime-switching 模型与基于聚类/HMM 的方法，用来判断当前市场处于哪种制度，进而切换风险预算。这一环与本文反复出现的“非平稳”约束直接对应——\textbf{制度识别做得好，等于给所有上游模型加了一个“当前学到的规律还成立吗”的开关。}

\subsection{环节耦合：本文最想强调的一点}

把五个环节放在一起看，会发现它们的优化目标互相冲突，而 AI 若只在单环节优化，往往会把冲突放大到别的环节：

\begin{table}[htbp]
\centering
\caption{交易环节之间的典型冲突}
\label{tab:conflict}
\begin{tabular}{@{}>{\raggedright\arraybackslash}p{3.2cm}>{\raggedright\arraybackslash}p{4.4cm}>{\raggedright\arraybackslash}p{7.6cm}@{}}
\toprule
耦合对 & 冲突 & 后果 \\
\midrule
信号 vs 执行 & 追求信号锐度会提高换手 & 成本吞噬 alpha；低信噪比下尤甚 \\
信号 vs 风控 & 追求收益会集中在高暴露 & 回撤放大，尾部风险上升 \\
组合 vs 执行 & 连续权重在小资金/高成本下不可行 & 理论组合与可实现组合差距大 \\
风控 vs 收益 & 波动率目标化会削峰 & 长期夏普可能改善但牛市落后 \\
制度识别 vs 上游 & 制度切换时上游模型失效 & 若不联动，模型在错误制度上继续下注 \\
\bottomrule
\end{tabular}
\end{table}

这张表想说的是：\textbf{端到端学习之所以被寄予厚望，正是因为它有机会在同一个目标下协调这些冲突；但它之所以难，也正因为它把这些冲突全部塞进了一个更难训练、更难诊断的黑箱。}对落地而言，一个务实的态度是：\textbf{理解耦合、显式地在上游为目标函数留出下游的代价，而不是把全部希望寄托在端到端上}。

本章要记住的一句话是：\textbf{一笔交易的成败，往往不在你能不能预测对方向，而在你把预测翻译成仓位、把仓位执行成成交、再把暴露管到明天的这三段衔接里——而这三段，恰恰是论文最不愿意细写的地方。}
